Whether we mean planning, justifying, and implementing a program for a non-profit or evaluating an already existing policy, causal inference to estimate the efficacy of policy is critical to decision making in government, business, and beyond. In general, policy analysts cannot conduct randomized trials, given their cost, ethical status, or infeasibility, and must therefore resort to carefully tuned econometric models to estimate counterfactuals, or a hypothetical scenario of how an outcome would have looked absent a given treatment or policy. Synthetic control methods (SCM) and difference-in-differences (DID) are among the most popular methods of estimating counterfactuals for policy analysts, and have seen much development in application and theory in recent years. In particular, some of their growth can be attributed to their extension via the integration of machine-learning methods. As such, my dissertation consists of three papers under a single theme: causal policy analysis in the high-dimensional setting, where \emph{high-dimensional} refers to where we have many more control units relative to pre-treatment periods.
